\chapter{콘텐츠 (읽기 전용)}

\begin{summarybox}{한눈에}
본사가 만든 모든 학습 콘텐츠 (일일 퀴즈·일일 독해·농장별·프로 모드·문법·작문 등) 를 \textbf{읽기 전용}으로 둘러봅니다. 학생에게 어떤 콘텐츠를 배정할지 미리 살펴볼 때, 학습 계획표 셀에 배정할 콘텐츠를 결정할 때 사용해요.
\end{summarybox}

\kfShot{06-content}{콘텐츠 목록 — 본사 모든 콘텐츠 (수정 불가)}

\section*{화면 구조}

\begin{screenbox}{콘텐츠 카탈로그}
\noindent
\begin{tabular}{@{}p{3cm}p{2cm}p{2cm}p{2.5cm}p{2cm}p{2cm}@{}}
\toprule
\textbf{제목} & \textbf{레벨} & \textbf{영역} & \textbf{모듈 키} & \textbf{문항} & \textbf{미리보기} \\
\midrule
일일 퀴즈 1일차 & 소쉬르1 & 비문학 & daily\_quiz & 10 & \inacttab{미리보기} \\
독해 훈련 5 & 러셀2 & 문학 & reading\_training & 5 & \inacttab{미리보기} \\
\bottomrule
\end{tabular}

\medskip
\textbf{우측 상단}: ``읽기 전용'' 배지 — 본사 관리자가 만들고 검수한 자료를 그대로 사용
\end{screenbox}

\section*{여기서 할 수 있는 일}
\begin{itemize}[leftmargin=1.2em]
  \item 영역(비문학·문학·문법·화법·작문·매체·논리·배경) · 레벨 · 모듈 키로 \textbf{검색·필터}
  \item \textbf{미리보기 클릭} → 해당 콘텐츠 학습 모듈을 그대로 시연 (학생이 보는 것과 동일)
  \item 검색 결과를 학습 계획표 셀 배정 시 그대로 활용
\end{itemize}

\section*{자주 쓰는 흐름}
\begin{enumerate}
  \item 학습 계획표 만들기 전에 이 화면에서 \textbf{어떤 콘텐츠를 쓸지} 미리 봅니다
  \item 영역·레벨·모듈 키 필터로 후보 추리기
  \item 미리보기로 학생 시점 확인
  \item 학습 계획표 셀 배정 모달의 ``콘텐츠 검색'' 으로 동일 카탈로그 사용
\end{enumerate}

\begin{tipbox}{알아두면 좋은 점}
\begin{itemize}[leftmargin=1.2em]
  \item 기관 관리자에게는 \textbf{[표준 양식][업로드][편집][삭제][일괄 삭제]} 버튼이 모두 숨김 — 본사가 운영하는 콘텐츠 풀에 직접 변경할 수 없음
  \item 화면 우측 상단의 ``(읽기 전용)'' 라벨이 권한 분기를 표시
  \item 본인 기관 안에서 만들고 싶은 학습은 \textbf{내용 숙지} 메뉴에서 별도로 가능 (Chapter 7)
\end{itemize}
\end{tipbox}
